% o texto deve começar imediatamente na próxima linha a fim de evitar identação
Este trabalho apresenta o desenvolvimento de um sistema embarcado de baixo
custo para o monitoramento de grandezas elétricas, com foco na avaliação
da qualidade da energia fornecida em instalações residenciais. O sistema
proposto é construído em torno do microcontrolador Raspberry Pi Pico 2 W
(RP2350), responsável pela aquisição, processamento e registro dos dados
de tensão e corrente. A aquisição de corrente é realizada por meio de
sensores SCT013-000, um por fase, enquanto a tensão é obtida com um
sensor ZMPT101B, com a seleção dos canais realizada por um multiplexador
CD74HC4067. Os dados processados são armazenados em cartão microSD no
formato CSV e exibidos localmente em um display OLED SSD1306, além de
sincronizados temporalmente por meio do protocolo NTP. Até o presente
estágio do trabalho, foram concluídas as etapas de condicionamento
analógico dos sinais, desenvolvimento do firmware embarcado, calibração
dos sensores de tensão e corrente e validação inicial do sistema por
meio de ensaios em regime permanente e sob condição de falta de fase. Os
resultados parciais obtidos indicam que a arquitetura desenvolvida é
capaz de realizar a aquisição contínua das grandezas elétricas e
registrar corretamente a ocorrência de distúrbios reais na rede.
Contudo, restam a implementação do cálculo de potência ativa, reativa e
fator de potência, o desenvolvimento de uma interface web para
monitoramento remoto, a validação do sistema com bateria de \textit{backup}
para garantir a continuidade do registro em caso de interrupção do
fornecimento, o desenvolvimento de um invólucro para acondicionamento
físico do sistema e, por fim, a implementação de um modelo de
aprendizado de máquina para previsão do consumo energético a partir dos
dados coletados. Conclui-se que a solução proposta é viável técnica e
economicamente, apresentando potencial para aplicações de monitoramento
residencial de baixo custo e identificação de eventos que possam
comprometer a qualidade da energia entregue ao consumidor.


\palavraschave{Qualidade de Energia Elétrica, Sistema Embarcado, Raspberry Pi Pico 2 W,
Aprendizado de Máquina, Internet das Coisas (IoT).}
% ATENÇÃO! Obrigatório um espaço entre a última linha do resumo e o comando \palavraschave{...}
